\section{From the $A_\infty$–chiral structure to a Poisson vertex algebra on cohomology}

\subsection{Setup and standing hypotheses}

Let $A$ denote the cochain complex of bulk quantum local operators of a $3$–dimensional holomorphic–topological field theory on $\R_t\times \C_z$, equipped with:
\begin{itemize}[leftmargin=2em]
 \item a BV–BRST differential $Q\colon A\to A$ of cohomological degree $+1$;
 \item an $A_\infty$–chiral structure $\{m_k\}_{k\ge1}$, with
 \begin{equation}
 % label removed: eq:A-infty-chiral
 m_k \colon A^{\otimes k} \longrightarrow A((\lambda_1))\cdots((\lambda_{k-1})) \qquad\text{of degree } 2-k,
 \end{equation}
 satisfying the (chiral) Stasheff relations, and with $m_1=Q$.
\end{itemize}

The analytic locality and renormalization hypotheses are as follows.

\begin{assumption}[Analytic hypotheses (H1)–(H4)]
% label removed: assump:H1-H4
The propagators and interaction densities obtained from the (renormalized) BV action of the holomorphic–topological theory on $\R\times \C$ are such that:
\begin{enumerate}[label=(H\arabic*),leftmargin=2em]
 \item All correlation kernels that define the $m_k$ are meromorphic in the holomorphic variables and rapidly decaying in the topological variable; 
 \item the operations $m_k$ are defined by weight forms with at most simple poles along the collision divisors $D_S \subset \FM_k(\C)$. Equivalently, the weight forms are logarithmic in the sense that they involve factors of the form $d\log(z_i - z_j) = dz_i/(z_i - z_j)$. These singularities are integrable in the BV sense (no new UV divergences arise in defining $m_k$);
 \item composition of operations is compatible with meromorphic factorization, i.e. factorization of kernels mirrors boundary decompositions of Fulton–MacPherson compactifications $\FM_k(\C)$;
 \item all renormalization–scheme ambiguities in defining ($m_k$) are $Q$–exact homotopies (hence immaterial at the cohomology level).
\end{enumerate}
\end{assumption}

\begin{remark}
Assumption~\ref{assump:H1-H4} is satisfied by the holomorphic--topological theories on $\R\times\C$ treated by the BV renormalization framework; it is the hypothesis under which the configuration--space proofs below are valid.
\end{remark}

Let $H:=H^\bullet(A,Q)$ denote the cohomology of $(A,Q)$. Our aim is to extract from $\{m_k\}$ a \emph{Poisson vertex algebra} structure on $H$, i.e.:
\begin{itemize}[leftmargin=2em]
 \item a commutative, associative degree–$0$ product $\cdot\colon H\otimes H\to H$,
 \item a degree–$(-1)$ $\lambda$–bracket $\{\,\cdot {}_\lambda \cdot\,\}\colon H\otimes H\to H[\![\lambda]\!]$,
\end{itemize}
satisfying sesquilinearity, skewsymmetry, the Jacobi identity, and the Leibniz rule (vertex Poisson compatibility).

\subsection{Laurent decomposition and the regular/singular splitting of $m_2$}

The binary operation $m_2$ splits canonically into regular and singular parts in the spectral parameter.

\begin{definition}[Laurent projections]
% label removed: def:Laurent-projection
For a Laurent series $F(\lambda)\in A((\lambda))$, write
\[
F(\lambda) = F^{\ge0}(\lambda) + F^{<0}(\lambda),
\]
where $F^{\ge0}(\lambda)\in A[\![\lambda]\!]$ is the regular (non–negative power) part and $F^{<0}(\lambda)\in \lambda^{-1}A[\![\lambda^{-1}]\!]$ is the singular (negative power) part. We define
\[
m_2^{\mathrm{reg}}(a,b) := \Big(m_2(a,b)\Big)^{\ge0}, 
\qquad 
m_2^{\mathrm{sing}}(a,b) := \Big(m_2(a,b)\Big)^{<0}.
\]
\end{definition}

Both $m_2^{\mathrm{reg}}$ and $m_2^{\mathrm{sing}}$ have degree $-1$ as cochains. We use graded symmetrization/antisymmetrization:
\[
\Sym(X,Y):=\frac12\big( X + (-1)^{\degree{X}\degree{Y}}Y\big),
\qquad
\ASym(X,Y):=\frac12\big( X - (-1)^{\degree{X}\degree{Y}}Y\big).
\]

\subsection{Definition of the product on cohomology}

\begin{definition}[Classical product on $H$]% label removed: def:product
For classes $[a],[b]\in H$, choose representatives $a,b\in A$ and define
\begin{equation}% label removed: eq:product-def
[a]\cdot [b] \;:=\; \big[\, \Sym\big(m_2^{\mathrm{reg}}(a,b),\, m_2^{\mathrm{reg}}(b,a)\big)\,\big] \;\in H.
\end{equation}
\end{definition}

\begin{remark}[Degree accounting]% label removed: rem:degree-shift
The operation $m_2$ has cohomological degree~$-1$. The spectral parameter $\lambda$ carries degree~$+1$, so the coefficient of $\lambda^n$ in $m_2(a,b)$ has degree $|a|+|b|-1-n$. The product is defined via the $\lambda = 0$ coefficient, hence has degree $|a|+|b|-1$ as an operation: it is a degree~$-1$ product on~$A$. In the $(-1)$-shifted setting, this is a degree-$0$ product on the desuspension $s^{-1}A$, matching the usual convention. Similarly, the $\lambda$-bracket $\{a_\lambda b\} = [m_2(a,b)]_{\mathrm{sing}}$ has degree~$-1$ as a map $A \otimes A \to A((\lambda))$.
\end{remark}

\begin{lemma}[Well–definedness; \ClaimStatusProvedHere]% label removed: lem:product-welldefined
The operation \eqref{eq:product-def} is independent of the choice of representatives of $[a]$ and $[b]$.
\end{lemma}

\begin{proof}
Suppose $a\mapsto a+Qx$; using the $A_\infty$ identity at arity $n=2$,
\[
m_1(m_2(a,b)) + m_2(m_1(a),b) + (-1)^{\degree{a}} m_2(a, m_1(b)) \;=\; 0,
\]
where the sign $(-1)^{\degree{a}}$ on the third term is the standard Koszul sign from passing $m_1$ past $a$. This uses the fact that $Q = m_1$ commutes with the Laurent projection, since $Q$ acts on the field-algebra coefficients and does not involve the spectral parameter $\lambda$. Projecting to the regular part, we see that the change in $\Sym(m_2^{\mathrm{reg}}(a,b),m_2^{\mathrm{reg}}(b,a))$ is $Q$–exact. The same holds for $b\mapsto b+Qy$. Thus the class \eqref{eq:product-def} is independent of representatives.
\end{proof}

\begin{lemma}[Graded commutativity on $H$; \ClaimStatusProvedHere]% label removed: lem:commutativity
For all $[a],[b]\in H$ we have $[a]\cdot [b] = (-1)^{\degree{a}\degree{b}} [b]\cdot [a]$.
\end{lemma}

\begin{proof}
Immediate from the graded symmetrization in \eqref{eq:product-def}.
\end{proof}

\begin{proposition}[Associativity on $H$; \ClaimStatusProvedHere]% label removed: prop:associativity
The product $\cdot$ is associative on $H$.
\end{proposition}

\begin{proof}
Consider the chain–level associator 
\[
\alpha(a,b,c) := m_2^{\mathrm{reg}}(m_2^{\mathrm{reg}}(a,b),c) - m_2^{\mathrm{reg}}(a, m_2^{\mathrm{reg}}(b,c))\, .
\]
The $A_\infty$ identity at arity $3$ implies $Qm_3(a,b,c)$ equals a signed sum of such associators up to $Q$–exact terms. Projecting to the regular part and passing to cohomology, we obtain $[\alpha(a,b,c)]=0$ in $H$. Carefully tracking the graded symmetrizations in the outer and inner arguments yields associativity for $\cdot$.
\end{proof}
\subsection{Definition of the $\lambda$–bracket on cohomology}
\begin{definition}[Vertex Poisson bracket]% label removed: def:lambda-bracket
For $[a],[b]\in H$, define
\begin{equation}% label removed: eq:lambda-bracket
\{\, [a] {}_\lambda [b] \,\} \;:=\; \Big[\, \tfrac{1}{2}\big( m_2^{\mathrm{sing}}(a,b;\lambda) \;-\; (-1)^{(\degree{a}+1)(\degree{b}+1)}\, e^{-(\lambda+\partial)\partial_\lambda}\, m_2^{\mathrm{sing}}(b,a;\lambda)\big)\,\Big] \;\in H[\![\lambda]\!].
\end{equation}
Here the operator $e^{-(\lambda+\partial)\partial_\lambda}$ implements the vertex algebra skew-symmetry transformation $\lambda\mapsto -\lambda-\partial$, which is \emph{not} the same as naive antisymmetrization for $\lambda$-dependent series (cf.\ \eqref{eq:bracket-def} in axioms).
\end{definition}
\begin{lemma}[Well–definedness of the $\lambda$–bracket; \ClaimStatusProvedHere]% label removed: lem:bracket-welldefined
The operation \eqref{eq:lambda-bracket} is independent of the choice of representatives of $[a],[b]$.
\end{lemma}
\begin{proof}
As in Lemma~\ref{lem:product-welldefined}, $a\mapsto a+Qx$ shifts $m_2(a,b)$ by $Q$–exact terms; projecting to the singular part and antisymmetrizing still yields a $Q$–exact variation. Hence the cohomology class in \eqref{eq:lambda-bracket} is independent of representatives.
\end{proof}

\subsection{Sesquilinearity and skewsymmetry}

\begin{proposition}[Sesquilinearity and skewsymmetry; \ClaimStatusProvedHere]% label removed: prop:sesqui-skew
The bracket $\{\cdot {}_\lambda \cdot\}$ on $H$ satisfies:
\begin{enumerate}[label=(\alph*),leftmargin=2em]
 \item skewsymmetry: $\{\, [a] {}_\lambda [b] \,\} = - (-1)^{(\degree{a}+1)(\degree{b}+1)} \{\, [b] {}_{-\lambda-\partial} [a] \,\}$;
 \item sesquilinearity: $\{\,\partial [a] {}_\lambda [b] \,\} = -\lambda \{\, [a] {}_\lambda [b] \,\}$ and $\{\, [a] {}_\lambda \partial [b] \,\} = (\lambda+\partial)\{\, [a] {}_\lambda [b] \,\}$.
\end{enumerate}
Here $\partial$ denotes the operator implementing translation along the holomorphic coordinate (defined earlier in the paper in the BV/factorization formalism) and acts on $H$ by the induced derivation.
\end{proposition}
\begin{proof}
Sesquilinearity on cohomology follows from the chain-level sesquilinearity axioms (\eqref{eq:sesqui-left} and \eqref{eq:sesqui-right}): since both $\partial$ and the spectral parameters $\lambda_i$ commute with $Q$, the sesquilinearity relations descend to cohomology.

Skew-symmetry follows from the definition of the bracket via antisymmetrization with the exponential shift operator $e^{-(\lambda+\partial)\partial_\lambda}$, which by construction imposes $\{\, [a] {}_\lambda [b] \,\} = -(-1)^{(\degree{a}+1)(\degree{b}+1)}\{\, [b] {}_{-\lambda-\partial} [a] \,\}$.
\end{proof}
\subsection{Jacobi identity and derivation (Leibniz) property}
\begin{theorem}[Jacobi identity; \ClaimStatusProvedHere]% label removed: thm:Jacobi
For all $[a],[b],[c]\in H$ we have in $H[\![\lambda,\mu]\!]$:
\[
\{\, [a] {}_\lambda \{\, [b] {}_\mu [c] \,\} \,\} \;-\; (-1)^{(\degree{a}+1)(\degree{b}+1)} \{\, [b] {}_\mu \{\, [a] {}_\lambda [c] \,\} \,\} \;=\; \{\, \{\, [a] {}_\lambda [b] \,\} {}_{\lambda+\mu} [c] \,\}.
\]
\end{theorem}
\begin{proof}
At the chain level, the Jacobi identity is encoded by the $A_\infty$ identity at arity $3$, projected to the \emph{singular} part of three–point functions. The four boundary strata of $\FM_3(\C)$ are $D_{12}$, $D_{13}$, $D_{23}$, and $D_{123}$. The stratum $D_{123}$ corresponds to the total collision of all three points. Its contribution is $Q$-exact (it corresponds to $m_1(m_3(a,b,c))$ in the $A_\infty$ identity) and vanishes on cohomology classes. The three pairwise collision strata $D_{12}$, $D_{13}$, $D_{23}$ contribute the three terms of the Jacobi identity. Their cancellation follows from the classical Arnold relation among logarithmic 1-forms:
\[
 d\log(z_1-z_2) \wedge d\log(z_2-z_3) \;+\; d\log(z_2-z_3) \wedge d\log(z_3-z_1) \;+\; d\log(z_3-z_1) \wedge d\log(z_1-z_2) \;=\; 0.
\]
Passing to cohomology and using Lemma~\ref{lem:bracket-welldefined} yields the result.
\end{proof}
\begin{theorem}[Leibniz rule; \ClaimStatusProvedHere]% label removed: thm:Leibniz
For all $[a],[b],[c]\in H$ we have:
\[
\{\, [a] {}_\lambda ([b]\cdot [c]) \,\} = \{\, [a] {}_\lambda [b] \,\}\cdot [c] + (-1)^{(\degree{a}+1)\degree{b}} \,[b]\cdot \{\, [a] {}_\lambda [c] \,\}.
\]
\end{theorem}
\begin{proof}
At the chain level this is the statement that the \emph{singular} part of $m_2$ acts as a bi–derivation with respect to the \emph{regular} part. Project the $n=3$ $A_\infty$ identity to the case where the outer operation contributes the regular part (product) and the inner operation contributes the singular part (bracket). Tracking the spectral parameter substitution rule from \eqref{eq:ainfty-relation-raw}, the inner bracket at position $s$ carries spectral parameter $\Lambda_{[s+1,\ldots,s+j]}$. For $j=2$ (inner bracket) and $k=2$ (outer product), the identity gives precisely the Leibniz compatibility between the singular and regular parts of $m_2$, with the precise degree signs being the standard Koszul signs induced by the degrees of $a,b,c$. Passing to cohomology yields the stated Leibniz rule.
\end{proof}

\begin{corollary}[PVA structure; \ClaimStatusProvedHere]% label removed: cor:PVA-structure
The pair $(\cdot,\,\{\,\cdot {}_\lambda \cdot\,\})$ defines a Poisson vertex algebra on $H$ of degree shift $-1$: the product $\cdot$ is graded commutative and associative, the $\lambda$–bracket is of degree $-1$ and satisfies sesquilinearity, skewsymmetry, the Jacobi identity, and acts by bi–derivations with respect to $\cdot$.
\end{corollary}

\subsection{Configuration space model and residue calculus (geometric proof outline)}

We record the geometric incarnation of the previous algebraic arguments. Denote by $\Conf_k(\C)$ the configuration space of $k$ ordered distinct points in $\C$ and by $\FM_k(\C)$ its Fulton–MacPherson compactification. As usual, the boundary strata of $\FM_k(\C)$ parametrize collisions of subsets of points with canonical sphere (direction) data. The $A_\infty$ operations $m_k$ are given by integrals of logarithmic (meromorphic) forms over $\FM_k(\C)$ built from propagators and interaction vertices; by (H1)–(H4), the only singularities are logarithmic along the exceptional divisors.

\begin{definition}[Regular/singular projections by pole order]% label removed: def:pole-projections
For the two–point operation $m_2$ the regular/singular decomposition of Definition~\ref{def:Laurent-projection} coincides with projection onto the smooth/logarithmic parts of the kernel with respect to the sole boundary divisor $D_{12}\subset \FM_2(\C)$. This identification holds because the spectral parameter $\lambda$ is conjugate to the FM coordinate $\varepsilon$ at the collision divisor $D_{12}$: near $D_{12}$, the propagator kernel takes the form $\sum_n a_{(n)}b \cdot \varepsilon^{-n-1}$, so the Laurent expansion in $\lambda^{-1}$ coincides with the expansion in powers of $\varepsilon^{-1}$ near the boundary.
\end{definition}

\begin{proposition}[Arnold relations and Jacobi; \ClaimStatusProvedHere]% label removed: prop:AOS-implies-Jacobi
The Arnold relation for the arrangement of diagonals $\{z_i = z_j\}$ in $\C^k$ states that for any triple $i<j<l$, the logarithmic 2-form $d\log(z_i - z_j) \wedge d\log(z_j - z_l) + \text{cyclic} = 0$. For $k=3$, this is the unique Arnold relation and it directly yields the Jacobi identity for the $\lambda$-bracket defined by residues along the collision divisors.
\end{proposition}

\begin{proof}
For $k=3$, the arrangement $\{z_1=z_2,\, z_2=z_3,\, z_1=z_3\}$ admits the single Arnold relation
\[
 d\log(z_1-z_2)\wedge d\log(z_2-z_3) \;+\; d\log(z_2-z_3)\wedge d\log(z_3-z_1) \;+\; d\log(z_3-z_1)\wedge d\log(z_1-z_2) \;=\; 0.
\]
Pushing forward along the forgetful maps to the $\lambda$–variables reproduces the stated Jacobi identity in vertex notation.
\end{proof}

\begin{remark}
The degree shifts (the $-1$ in the bracket) originate from the fact that $m_2$ has cohomological degree $-1$ at the chain level: geometrically this is the degree of the current defined by the logarithmic kernel (one complex logarithmic pole contributes cohomological degree $-1$ after BV parity conventions).
\end{remark}


